\section*{Chapitre \ref{ch:b2:maxwell-equations} --- Les équations de Maxwell}

\begin{solution}{exo:b2:maxwell-equations:1}
(a) $\operatorname{div} = 3$, $\operatorname{\vect{curl}} = \vect 0$~: le \omterm{def:b2:maxwell-equations:operators}{gradient} de $r^2/2$.
(b) $0$ et $(0, 0, 2)$~: un \omterm{def:b2:maxwell-equations:operators}{rotationnel} (de $-\tfrac12(x^2 + y^2)\vect e_z$), non un
\omterm{def:b2:maxwell-equations:operators}{gradient}. (c) $0$ et $\vect 0$~: le \omterm{def:b2:maxwell-equations:operators}{gradient} de $xyz$ (et aussi un \omterm{def:b2:maxwell-equations:operators}{rotationnel}). (d)
$\frac1{r^2}\partial_r(r^2/r^2) = 0$, \omterm{def:b2:maxwell-equations:operators}{rotationnel} $\vect 0$~: le \omterm{def:b2:maxwell-equations:operators}{gradient} de $-1/r$. (e)
$0$ et $(0, -2x, 0)$~: un \omterm{def:b2:maxwell-equations:operators}{rotationnel}, non un \omterm{def:b2:maxwell-equations:operators}{gradient}.
\end{solution}

\begin{solution}{exo:b2:maxwell-equations:2}
(a) Gauss et Faraday~: les deux membres en \unit{V/m^2} (une \omterm{def:b2:charges-currents-conduction:densities}{densité de charge} divisée par
$\varepsilon_0$, un \omterm{def:b2:maxwell-equations:operators}{gradient} de champ, un champ divisé par un temps)~; Ampère~: les deux membres en
\unit{T/m}.
(b) $1/\sqrt{8.85 \times 10^{-12} \times 1.257 \times 10^{-6}} = \qty{3.00e8}{m/s}$. (c) Rapport
$\varepsilon_0\omega/\gamma$~: cuivre $5 \times 10^{-20}$~; eau de mer à \qty{1}{GHz} $1 \times 10^{-2}$
(encore un conducteur, de justesse)~; verre à \qty{50}{Hz} $3 \times 10^3$~: un diélectrique,
le \omterm{thm:b2:maxwell-equations:maxwell}{courant de déplacement} domine.
\end{solution}

\begin{solution}{exo:b2:maxwell-equations:3}
$\dd E/\dd t = I/\varepsilon_0\pi R^2 = \qty{1.8e12}{V.m^{-1}.s^{-1}}$~; $j_d = I/\pi R^2 =
\qty{16}{A/m^2}$~; $B(\qty{5}{cm}) = \mu_0Ir/2\pi R^2 = \qty{0.5}{\micro T}$~; $B(\qty{20}{cm}) =
\mu_0I/2\pi r = \qty{0.5}{\micro T}$ --- exactement le champ du fil.
\end{solution}

\begin{solution}{exo:b2:maxwell-equations:4}
(a) $\rho = \varepsilon_0k$. (b) $\frac1{r^2}\partial_r(\rho_0r^3/3\varepsilon_0) = \rho_0/\varepsilon_0$~;
à l'extérieur, Gauss~: $E = \rho_0R^3/3\varepsilon_0r^2$. (c) $\operatorname{div}\vect B = 0$~: oui~;
$\operatorname{\vect{curl}}\vect B = -(B_0/a)\vect e_z$~: $\vect j = -(B_0/\mu_0a)\vect e_z$, une
nappe de courant uniforme d'épaisseur selon $y$.
\end{solution}

\begin{solution}{exo:b2:maxwell-equations:5}
(a) $\vect A = \tfrac12(B_yz - B_zy, B_zx - B_xz, B_xy - B_yx)$~: $(\operatorname{\vect{curl}}
\vect A)_x = \tfrac12(B_x + B_x) = B_x$, etc.~; $\operatorname{div}\vect A = 0$. (b)
$\vect A' - \vect A = \tfrac12B_0(y, x, 0) = \operatorname{\vect{grad}}(\tfrac12B_0xy)$. (c)
$\frac1r\partial_r(Br^2/2) = B$~; $\frac1r\partial_r(BR^2/2) = 0$~; les deux valent $BR/2$ en
$R$. (d) $2\pi r\cdot BR^2/2r = \pi R^2B = \Phi$~: un contour extérieur au solénoïde,
là où $\vect B = \vect 0$, ressent tout de même $e = -\dd\Phi/\dd t$ --- par $\vect E =
-\partial_t\vect A$, qui ne s'y annule pas.
\end{solution}

\begin{solution}{exo:b2:maxwell-equations:6}
(a) $E = \sigma/\varepsilon_0 = \qty{1.1e5}{V/m}$, normal. (b) $\sigma/\varepsilon_0$ entre les plans, nul
à l'extérieur (chaque saut ajoute $\pm\sigma/\varepsilon_0$). (c) $j_s = nI = \qty{1e4}{A/m}$, saut
$\mu_0j_s = \qty{12.6}{mT} = B$ à l'intérieur. (d) Un contour extérieur au tore n'enlace
aucun courant net, donc $B = 0$ en ce lieu~; à l'intérieur, Ampère donne $\mu_0NI/2\pi r$,
et le saut à la traversée de l'enroulement, $\mu_0j_s = \mu_0NI/2\pi r$, concorde.
\end{solution}

\begin{solution}{exo:b2:maxwell-equations:7}
(a) $V'' = -\rho_0/\varepsilon_0$~: $V = -\rho_0x^2/2\varepsilon_0 + (U/d + \rho_0d/2\varepsilon_0)x$. (b)
$E = \rho_0x/\varepsilon_0 - U/d - \rho_0d/2\varepsilon_0$, nul en $x_0 = d/2 + \varepsilon_0U/\rho_0d$
--- près du milieu pour une grande charge d'espace. (c) $\sigma(0) = \varepsilon_0E(0^+) =
-\varepsilon_0U/d - \rho_0d/2$, $\sigma(d) = -\varepsilon_0E(d^-) = \varepsilon_0U/d - \rho_0d/2$~: au total
$-\rho_0d$, qui neutralise la charge d'espace. (d) $\rho_0 = 0$~: $V$ linéaire, $E$
uniforme~; $U = 0$~: une parabole culminant en $d/2$.
\end{solution}

\begin{solution}{exo:b2:maxwell-equations:8}
(a) \qty{6000}{km}, \qty{300}{m}, \qty{3}{m}, \qty{12.5}{cm}. (b) $f < c/10L$~:
\qty{300}{MHz}~; \qty{3}{GHz}~; \qty{30}{Hz}. (c) $(2\pi \times 10^6 \times 0.05/3 \times 10^8)^2/4 =
3 \times 10^{-7}$. (d) Il ne l'est pas, à la rigueur~: à \qty{50}{Hz} un réseau de \qty{1000}{km} vaut
un sixième de longueur d'onde, les phases y diffèrent, et les longues lignes sont
traitées comme des \omterm{prop:b2:dispersion-wave-packets:coax}{lignes de transmission} --- chaque circuit local, lui, est
quasi stationnaire.
\end{solution}

\begin{solution}{exo:b2:maxwell-equations:9}
$\operatorname{div}\operatorname{\vect{curl}}\vect B = 0$ imposerait $\operatorname{div}\vect j =
0$, c'est-à-dire $\partial_t\rho = 0$ partout. Dans l'armature (volume $\pi R^2e$) le
courant $I$ entre et rien ne sort~: $\langle\operatorname{div}\vect j\rangle =
-I/\pi R^2e = \qty{-1.6e5}{A/m^3}$, $\partial_t\rho = \qty{+1.6e5}{C.m^{-3}.s^{-1}}$.
\end{solution}

\begin{solution}{exo:b2:maxwell-equations:10}
(a) $2\pi rE_\theta = -\pi r^2\dot B$~: $E_\theta = -r\dot B/2$ à l'intérieur, $-R^2\dot B/2r$
à l'extérieur. (b) Ses lignes sont des cercles fermés~: ce n'est pas un \omterm{def:b2:maxwell-equations:operators}{gradient}, il
entraîne les charges le long d'un contour --- un champ électromoteur. (c) $e\cdot2\pi r_0
|E_\theta| = e\pi r_0^2\dot B$~: $\pi \times 0.25 \times 10 = \qty{7.9}{eV}$ par tour~; $2.5 \times 10^6$
tours. (d) $p = eB(r_0)r_0$ donne $\dot p = er_0\dot B(r_0)$, tandis que le champ
induit donne $\dot p = e|E_\theta| = e\dot\Phi/2\pi r_0$~: $B(r_0) = \Phi/2\pi r_0^2 =
\tfrac12\langle B\rangle$.
\end{solution}

\begin{solution}{exo:b2:maxwell-equations:11}
(a) $\operatorname{\vect{curl}}\operatorname{\vect{curl}}\vect E = -\Delta\vect E$ (puisque $\operatorname{div}
\vect E = 0$) $= -\partial_t\operatorname{\vect{curl}}\vect B = -\mu_0\varepsilon_0\partial_t^2\vect E$. (b)
De même avec les rôles échangés. (c) $c = 1/\sqrt{\mu_0\varepsilon_0} = \qty{3.0e8}{m/s}$.
(d) Weber et Kohlrausch avaient mesuré le rapport des unités électrostatique
et électromagnétique de charge, $3.1 \times 10^8$ m/s~; Fizeau avait
mesuré la lumière à $3.15 \times 10^8$ m/s~: «~nous pouvons difficilement éviter
d'en conclure que la lumière consiste en ondulations transversales du
même milieu~».
\end{solution}

\begin{solution}{exo:b2:maxwell-equations:12}
(a) $\operatorname{\vect{curl}}\operatorname{\vect{curl}}\vect B = -\Delta\vect B = \mu_0\gamma\operatorname{\vect{curl}}
\vect E = -\mu_0\gamma\partial_t\vect B$~: $\Delta\vect B = \mu_0\gamma\partial_t\vect B$, la forme de
l'équation de la chaleur avec la diffusivité $1/\mu_0\gamma$. (b) $\tau \sim \mu_0\gamma d^2$~: \qty{7.5}{ms}
pour la plaque~; $4\pi \times 10^{-7} \times 5 \times 10^5 \times 9 \times 10^{12} = \qty{6e12}{s}$, deux cent
mille ans, pour le noyau. (c) $\mu_0\gamma\delta^2 \sim 2/\omega$~: $\delta = \sqrt{2/\mu_0
\gamma\omega}$~: \qty{9}{mm} à \qty{50}{Hz}, \qty{65}{\micro m} à \qty{1}{MHz}. (d) Le champ
du noyau ne peut pas changer par diffusion en moins de $10^5$ ans~: la
variation séculaire observée et les inversions exigent le mouvement du
fluide conducteur --- une dynamo.
\end{solution}

\begin{solution}{pb:b2:maxwell-equations:1}
\textbf{1.} $Q = (I_0/\omega)\sin\omega t$, $E = Q/\varepsilon_0\pi R^2$~: amplitude $I_0/\omega\varepsilon_0
\pi R^2 = \qty{5.7e5}{V/m}$.

\textbf{2.} $j_d = \varepsilon_0\partial_tE = I(t)/\pi R^2 = 32\cos\omega t$ A/m$^2$~; son flux
à travers l'armature vaut $I(t)$.

\textbf{3.} $2\pi rB = \mu_0\varepsilon_0\pi r^2\partial_tE$~: $B = \mu_0I(t)r/2\pi R^2$~; en $R$~:
$\mu_0I_0/2\pi R = \qty{2}{\micro T}$.

\textbf{4.} Pour $r > R$ tout le \omterm{thm:b2:maxwell-equations:maxwell}{courant de déplacement} est enlacé~: $B =
\mu_0I/2\pi r$, le champ du fil~; continu en $R$.

\textbf{5.} $(\operatorname{\vect{curl}}\vect E_1)_\theta = -\partial_rE_1 = -\partial_tB_\theta$, donc $\partial_rE_1
= \mu_0\varepsilon_0(r/2)\ddot E$ et $E_1(r) - E_1(0) = \mu_0\varepsilon_0r^2\ddot E/4$~; avec
$\ddot E = -\omega^2E$~: $|E_1(R) - E_1(0)|/|E| = (\omega R/c)^2/4 = 1 \times 10^{-6}$.

\textbf{6.} Excellent. $10\%$ exigerait $(\omega R/c)^2 = 0.4$, $f \approx \qty{300}{MHz}$ ---
une cavité, et non plus un condensateur.

\textbf{7.} $W_e = \tfrac12\varepsilon_0E_0^2\pi R^2d = \qty{4.5e-5}{J}$~; $W_m = \int_0^R(\mu_0I_0r/
2\pi R^2)^2/2\mu_0\cdot2\pi rd\,\dd r = \mu_0I_0^2d/16\pi = \qty{2.5e-11}{J}$~; rapport $6 \times
10^{-7} \approx (\omega R/c)^2/8$~: le même petit paramètre.

\textbf{8.} $\oint\vect E\cdot\dd\vect l = 2\pi r_0E_\theta = -\dd\Phi/\dd t$.

\textbf{9.} $\dd p/\dd t = e|E_\theta| = (e/2\pi r_0)|\dd\Phi/\dd t|$~; par tour,
$2\pi r_0\cdot e|E_\theta| = e|\dd\Phi/\dd t|$.

\textbf{10.} $p = eB(r_0)r_0$ donne $\dot p = er_0\dot B(r_0)$~; en égalant avec 9~:
$B(r_0) = \Phi/2\pi r_0^2 = \tfrac12\langle B\rangle$.

\textbf{11.} Un champ uniforme a $B(r_0) = \langle B\rangle$, deux fois trop~:
la force magnétique croît plus vite que la quantité de mouvement et l'orbite
se resserre. Les pôles sont profilés pour que le champ soit fort près de
l'axe et deux fois plus faible sur l'orbite.

\textbf{12.} $p = eBr_0 = \qty{3.2e-20}{kg.m/s}$~; $E = pc = \qty{9.6e-12}{J} = \qty{60}{MeV}$~;
$\Phi = 2\pi r_0^2B = \qty{0.63}{Wb}$ en \qty{5}{ms}~: \qty{126}{eV} par tour, $4.8 \times
10^5$ tours, \qty{1500}{km} --- à $c$ en \qty{5}{ms}~: cohérent.

\textbf{13.} De l'alimentation de l'électroaimant, transmise aux électrons par
le terme de Faraday~: le flux variable crée le champ accélérateur.

\textbf{14.} Des dizaines de MeV d'électrons freinés dans une cible lourde donnent des
rayons X pénétrants~; l'accélérateur linéaire compact, de débit de dose plus élevé,
l'a remplacé.

\textbf{15.} $\operatorname{div}\operatorname{\vect{curl}}\vect A = 0$~; $\operatorname{\vect{curl}}(-\operatorname{\vect{grad}}
V - \partial_t\vect A) = -\partial_t\operatorname{\vect{curl}}\vect A = -\partial_t\vect B$.

\textbf{16.} $\frac1{r^2}(r^2V')' = -\rho_0/\varepsilon_0$~: $V = \rho_0(3a^2 - r^2)/6\varepsilon_0$
à l'intérieur, $\rho_0a^3/3\varepsilon_0r = Q/4\pi\varepsilon_0r$ à l'extérieur~; $E = \rho_0r/3\varepsilon_0$ et
$Q/4\pi\varepsilon_0r^2$.

\textbf{17.} $\tfrac12\int\rho V\,\dd\tau = 2\pi\rho_0^2(a^5/6 - a^5/30)/\varepsilon_0 = 4\pi\rho_0^2a^5/
15\varepsilon_0 = 3Q^2/20\pi\varepsilon_0a$~; uranium~: \qty{1.6e-10}{J} $\approx \qty{1000}{MeV}$ (la fission
libère les \qty{200}{MeV} dont deux boules plus petites coûtent moins cher).

\textbf{18.} $V = -E(t)z$ (à une constante près). Avec $\vect A = A_z(r, t)\vect e_z$,
$B_\theta = -\partial_rA_z = \mu_0Ir/2\pi R^2$ donne $A_z = -\mu_0I(t)r^2/4\pi R^2$.

\textbf{19.} $-\partial_tA_z = \mu_0\dot Ir^2/4\pi R^2 = \mu_0\varepsilon_0r^2\ddot E/4$ (puisque $\dot I =
\varepsilon_0\pi R^2\ddot E$)~: la correction de la question 5, à une constante près.

\textbf{20.} $\operatorname{div}\vect A = \partial_zA_z = 0$ parce que $A_z$ ne dépend que de
$r$~; n'importe quel $\vect A + \operatorname{\vect{grad}}\chi$.

\textbf{21.} $\sigma = \varepsilon_0E_0 = \qty{5}{\micro C/m^2}$~; dans le métal $\vect E =
\vect 0$, de sorte que le saut $\sigma/\varepsilon_0$ est exactement le champ de l'intervalle.

\textbf{22.} Le courant $I(t)$ alimente l'armature depuis l'axe~; la
part restant à livrer au-delà de $r$ vaut $I(1 - r^2/R^2)$, donc $j_s = I(1 -
r^2/R^2)/2\pi r$. Sur la face externe Ampère donne le champ du fil $\mu_0I/
2\pi r$~; sur la face interne $\mu_0Ir/2\pi R^2$~; la différence, $\mu_0I(1 -
r^2/R^2)/2\pi r = \mu_0j_s$, est le saut de la relation de passage.

\textbf{23.} Armatures~: $R/\lambda = 3 \times 10^{-4}$, \omterm{def:b2:maxwell-equations:arqs}{ARQS} électrique. Boucle~: $7 \times 10^{-4}$,
\omterm{def:b2:maxwell-equations:arqs}{ARQS} magnétique. À \qty{1}{GHz}~: $R/\lambda = 0.3$ --- une cavité et une antenne.

\textbf{24.} $\lambda = \qty{3000}{km}$ pour un anneau de l'ordre du mètre~: le \omterm{thm:b2:maxwell-equations:maxwell}{courant de
déplacement} est parfaitement négligeable~; le champ électrique essentiel est
celui de Faraday, que l'\omterm{def:b2:maxwell-equations:arqs}{ARQS} magnétique conserve en entier.

\textbf{25.} I~: Ampère--Maxwell (\omterm{thm:b2:maxwell-equations:maxwell}{courant de déplacement}), Faraday pour
la correction, densités d'énergie. II~: Faraday intégrale, force de Lorentz.
III~: les potentiels, Poisson (Gauss). IV~: les sauts de $\vect E$
et de $\vect B$ aux interfaces~; le rapport taille sur longueur d'onde pour le régime.
\end{solution}
